\documentclass[journal]{IEEEtran}

\usepackage{cite}
\usepackage{amsmath,amssymb,amsfonts}
\usepackage{graphicx}
\usepackage{textcomp}
\usepackage{xcolor}
\usepackage{booktabs}
\usepackage{multirow}
\usepackage[utf8]{inputenc}
\usepackage[T1]{fontenc}
\usepackage{microtype}
\usepackage{array}

\interdisplaylinepenalty=2500
\emergencystretch=2em
\hyphenation{Koop-man trajec-tory trajec-tories auto-regressive regu-lar-ization data-set data-sets non-over-lapping pre-train-ing Trans-former state-em-bedded delay-em-bed-ding}

\def\BibTeX{{\rm B\kern-.05em{\sc i\kern-.025em b}\kern-.08em
    T\kern-.1667em\lower.7ex\hbox{E}\kern-.125emX}}

\begin{document}

\title{Transformer-Koopman: Finite-History Lifting for Controlled Nonlinear Systems}

\author{Bolin Tao}

\markboth{IEEE Journal Draft}{Tao: Transformer-Koopman}

\maketitle

\begin{abstract}
Koopman modeling is attractive for controlled nonlinear systems because it seeks a lifted representation in which nonlinear dynamics become approximately linear. In practice, however, the quality of the resulting model depends critically on the lifting map. Existing learned Koopman models often construct observables from the instantaneous state alone, which can discard short-term temporal information useful for identifying accurate linear latent dynamics. This paper proposes Transformer-Koopman, a temporal lifting framework that learns Koopman observables from finite history windows by combining non-overlapping multivariate temporal patches with a Transformer encoder. The lifted state evolves according to linear latent dynamics controlled by the input, and the associated system matrices are identified through a hybrid pipeline that combines encoder pretraining with Extended Dynamic Mode Decomposition (EDMD) style regression. Open-loop experiments on a 6-state manipulator dataset show substantially lower rollout error than traditional EDMD and multi-layer perceptron (MLP) Koopman baselines, while ablation studies on a second 2-state dataset demonstrate that history length, patch design, and moderate regularization materially affect model quality. The resulting representation remains compact, interpretable, and naturally compatible with subsequent tracking-control studies.
\end{abstract}



\begin{IEEEkeywords}
Koopman operator, controlled nonlinear systems, temporal lifting, Transformer, data-driven modeling, trajectory prediction
\end{IEEEkeywords}

\section{Introduction}
\label{sec:introduction}

Learning predictive models for nonlinear controlled systems is a central problem in system identification, state estimation, and model-based control \cite{ljung1999system, vanoverschee1994n4sid, nelles2020nonlinear}. For control-oriented applications, model quality is not measured only by one-step fit, but also by whether the representation supports reliable multi-step prediction and a usable interface for downstream controller design. The Koopman framework is attractive in this regard because it seeks coordinates in which nonlinear state evolution is represented, at least approximately, by linear dynamics acting on observables \cite{koopman1931hamiltonian, koopman1932spectra, mezic2005spectral, rowley2009spectral, brunton2022modern}. When such a representation is available, prediction and control can be formulated in a linear latent space even though the physical system remains nonlinear \cite{otto2021koopman, proctor2018generalizing, korda2018linear}.

The principal obstacle is the lifting map. Classical Extended Dynamic Mode Decomposition (EDMD), kernel-based variants, and related transfer-operator approximations identify lifted linear predictors from predefined observable dictionaries \cite{williams2015data, williams2015kernel, klus2018data}. Their appeal lies in interpretability, spectral structure, and simple linear regression, but their performance depends strongly on the chosen observables. Learned liftings address this limitation by adapting the observable map directly from data \cite{takeishi2017koopman, yeung2019learning, lusch2018deep, morton2019deep, azencot2020forecasting}. However, many learned Koopman models still construct the lifted state from the instantaneous physical state alone. For controlled nonlinear systems, this can be restrictive when recent trajectory context carries information that is useful for producing an approximately linear latent evolution law.

This observation motivates a finite-history view of lifting. Delay coordinates, recurrent latent models, and linearly recurrent autoencoder architectures all suggest that recent temporal information can improve the representation of nonlinear dynamics \cite{arbabi2017ergodic, rank2019recurrent, otto2019linearly, forgione2021dynonet}. The question addressed in this paper is therefore not whether history can matter in principle, but how to use short histories without turning the lifted state into a full delay stack or an undifferentiated state augmentation. The proposed approach treats finite history as an input to a learned compression map: a short state-history window is encoded with non-overlapping temporal patches and a Transformer encoder, and the resulting temporal feature vector is appended to the current physical state. In this way, the model retains explicit state coordinates while allowing recent temporal context to shape the learned observables.

A second design requirement is compactness. The objective here is not generic sequence forecasting, but a control-compatible lifting map that can summarize recent trajectory information without introducing an unnecessarily heavy sequence model. Patch-wise tokenization provides a structured way to compress short local trajectory segments before global sequence processing, and recent compact-Transformer studies suggest that lightweight temporal encoders can already capture meaningful dynamical structure when tokenization is chosen carefully \cite{vaswani2017attention, nie2023time, zhou2022fedformer, liu2024itransformer, forootani2026compact, forootani2025synthetic}. In the present method, however, the Transformer is used only to construct observables from history. Prediction itself is produced by the identified linear latent model, so the paper should be read as a control-oriented model-identification study rather than as a forecasting architecture paper.

Based on this motivation, this paper proposes \emph{Transformer-Koopman}, a temporal lifting framework for controlled nonlinear systems. A history window is partitioned into non-overlapping multivariate temporal patches, embedded as tokens, processed by a Transformer encoder, and projected to a temporal feature vector. The lifted state is $z_k=[x_k;h_k]$, where the first $n$ coordinates are the physical state and the remaining coordinates are learned temporal features. This lifted state is propagated by the linear controlled model $z_{k+1}=Az_k+Bu_k$, and the physical state is recovered exactly by the fixed projection $x_k=Cz_k$ with $C=[I_n\;0]$. The resulting contribution is a control-compatible lifted predictor: the goal is to improve Koopman model identification while preserving a clean interface for subsequent controller design.

The contributions of this paper are summarized as follows.
\begin{enumerate}
    \item We propose a temporal lifting framework that constructs Koopman observables from finite history windows using non-overlapping temporal patches and a Transformer encoder, rather than from the current state alone.
    \item We construct a compact latent model in which the physical state is recovered by a fixed projection, preserving interpretability without requiring a learned nonlinear decoder.
    \item We provide completed open-loop experimental evidence through trajectory-prediction comparison and ablation analysis, and we summarize the control-oriented interface supported by the learned model.
\end{enumerate}

The claims of the paper are intentionally aligned with the available evidence. The completed results concern open-loop predictive modeling, not finished closed-loop validation. Accordingly, the main conclusion is that temporally informed lifting can yield a more accurate Koopman model, while the control section documents the model interface that would support later tracking studies.

The rest of the paper is organized as follows. Section~\ref{sec:related_work} reviews related work. Section~\ref{sec:problem} formulates the learning problem. Section~\ref{sec:methodology} presents the proposed temporal lifting method. Section~\ref{sec:control} discusses the control-oriented interface. Section~\ref{sec:analysis} discusses structural properties and limitations of the model. Section~\ref{sec:experiments} reports the current experiments. Section~\ref{sec:conclusion} concludes the paper.

\section{Related Work}
\label{sec:related_work}

\subsection{Koopman Modeling and Operator-Theoretic Foundations}

Koopman operator theory provides a linear description of nonlinear dynamics by acting on observables rather than directly on the state \cite{koopman1931hamiltonian, koopman1932spectra, mezic2005spectral, rowley2009spectral}. In data-driven settings, EDMD and related transfer-operator approximations provide finite-dimensional surrogates of this operator from trajectory data \cite{williams2015data, williams2015kernel, klus2018data}. These methods are appealing because they preserve linear latent dynamics and often admit interpretable spectral analysis. Their main limitation, however, is that performance depends strongly on the expressiveness of the chosen observables.

\subsection{Learned Liftings for Nonlinear Dynamics}

To reduce dependence on hand-crafted dictionaries, a large body of work has explored learned observable maps for Koopman modeling. Representative approaches include invariant-subspace learning, deep Koopman autoencoders, variational liftings, and consistency-based architectures \cite{takeishi2017koopman, yeung2019learning, lusch2018deep, morton2019deep, azencot2020forecasting, yuan2022variational}. Related linearly recurrent latent models and recurrent Koopman variants further emphasize that temporal structure can play an important role in learning linearizable representations of nonlinear dynamics \cite{otto2019linearly, rank2019recurrent}. More broadly, data-driven dynamical-system models such as DynoNet and sparse equation discovery highlight the value of combining structural priors with flexible function approximation \cite{forgione2021dynonet, brunton2016sindy}. The present work follows this line of thought, but specializes it to a controlled Koopman predictor whose latent structure remains explicitly anchored to the physical state.

\subsection{Koopman Models for Control and Prediction}

Koopman models are particularly attractive in control because linear latent dynamics can be combined with standard linear control tools. Dynamic Mode Decomposition with control, input-aware Koopman formulations, and Koopman-based MPC established this direction for controlled nonlinear systems \cite{proctor2016dmdc, proctor2018generalizing, brunton2016invariant, korda2018linear}. More recent work has studied learned Koopman representations for control, robust MPC, and closed-loop guarantees \cite{han2020deep, shi2022deep, strasser2026overview, duran2025control}. Taken together, this literature makes clear that controller performance depends heavily on model quality. For that reason, the main emphasis of this paper is on improving the learned lifting map itself, while keeping the resulting predictor structurally compatible with subsequent controller design.

\subsection{Transformers and Compact Temporal Encoders}

Transformers have become a standard sequence-modeling architecture because self-attention provides a flexible mechanism for capturing interactions across time \cite{vaswani2017attention}. In time-series modeling, representative architectures include efficient-attention and compact-attention variants \cite{zhou2021informer, wang2020linformer, katharopoulos2020transformers, choromanski2021performer}, decomposition-based models \cite{wu2021autoformer, zhou2022fedformer}, and simplified channel-time reformulations \cite{zeng2023transformers, liu2024itransformer}. Recent time-series encoders have also shown that patch-wise tokenization can reduce sequence length while preserving local temporal structure \cite{nie2023time}. This idea is particularly relevant for temporal lifting, because it offers a compact way to summarize recent trajectory segments before attention is applied.

For the present work, the most relevant references are not forecasting papers as such, but studies showing that compact Transformer variants can retain strong temporal representation capability under controlled synthetic settings \cite{forootani2026compact, forootani2025synthetic}. Closely related hybrid models have also begun to combine Transformer representations with Koopman-style latent dynamics \cite{forootani2026deepkoopformer, forootani2026spectral}. Our method differs from these forecasting-oriented studies in both objective and structure: it targets controlled nonlinear system identification, preserves the physical state explicitly in the latent vector, uses linear lifted dynamics identified by EDMD-style regression, and is positioned as a control-compatible Koopman predictor rather than a generic forecaster.

\section{Problem Formulation}
\label{sec:problem}

Consider a discrete-time controlled nonlinear system
\begin{equation}
    x_{k+1} = F(x_k,u_k),
    \label{eq:nonlinear_system}
\end{equation}
where $x_k \in \mathbb{R}^n$ is the state, $u_k \in \mathbb{R}^m$ is the control input, and $F(\cdot)$ is an unknown nonlinear transition map. We assume that full-state measurements and the corresponding control inputs are available from trajectory data. This assumption matches the datasets and learning pipeline considered in this study.

For a chosen history length $P$, define the history window
\begin{equation}
    X_k = [x_{k-P+1}^\top,\,x_{k-P+2}^\top,\,\ldots,\,x_k^\top]^\top \in \mathbb{R}^{P\times n}.
    \label{eq:history_window}
\end{equation}
From recorded trajectories, we form samples of the form
\begin{equation}
    \mathcal{D} = \{(X_k,u_k,x_{k+1})\}_{k=1}^{M},
    \label{eq:dataset}
\end{equation}
where only windows that remain within a single trajectory are retained. In the actual dataset construction, both the history window and the one-step prediction target are required to respect trajectory boundaries.

The modeling goal is to learn a finite-dimensional lifted representation in which the controlled nonlinear dynamics can be approximated by a linear latent system. To this end, we define a temporal lifting map
\begin{equation}
    z_k = g_\theta(X_k),
    \label{eq:temporal_lifting}
\end{equation}
where $z_k \in \mathbb{R}^d$ is a lifted state derived from a finite history window rather than from the instantaneous state alone. The target lifted dynamics are
\begin{equation}
    z_{k+1} = A z_k + B u_k + w_k,
    \label{eq:lifted_dynamics}
\end{equation}
where $A \in \mathbb{R}^{d\times d}$ and $B \in \mathbb{R}^{d\times m}$ are identified matrices and $w_k$ denotes the finite-dimensional approximation residual.

This formulation reflects the central hypothesis of the paper. If recent history carries information that is useful for constructing an approximately linear latent evolution law, then a finite-history lifting map can produce a more accurate Koopman model than an instantaneous-state encoder. At the same time, we do not simply treat the stacked history itself as the lifted state. Instead, we seek a compact lifted representation that preserves the current physical state explicitly while using the remaining latent coordinates to encode history-dependent information.

Two practical remarks are important. First, the present encoder uses \emph{state history only}; past controls are not included in the encoder input. This design keeps the lifting map simple and matches the training and data-construction procedure used here, and it should be understood as a modeling choice specific to the present formulation rather than as a general requirement of temporal lifting. Second, the primary objective of this paper is predictive model quality under autoregressive rollout, with control compatibility treated as a downstream structural requirement rather than as a completed closed-loop claim.

\section{Transformer-Based Temporal Lifting}
\label{sec:methodology}

This section presents the proposed model and its training procedure from the perspective of control-compatible model identification.

\subsection{Temporal Patch Tokenization}

The encoder takes the history window $X_k \in \mathbb{R}^{P\times n}$ as input. Assuming $p$ divides $P$, the window is partitioned into $N=P/p$ contiguous non-overlapping patches of temporal length $p$. Denote the $i$th patch by
\begin{equation}
    X_k^{(i)} \in \mathbb{R}^{p\times n}, \qquad i=1,\ldots,N.
\end{equation}
Each patch contains all $n$ state channels over $p$ consecutive time steps, so the tokenization is \emph{jointly multivariate} rather than channel-wise. After vectorization,
\begin{equation}
    q_k^{(i)} = \mathrm{vec}(X_k^{(i)}) \in \mathbb{R}^{pn},
\end{equation}
the corresponding token embedding is
\begin{equation}
    e_k^{(i)} = W_p q_k^{(i)} + b_p,
    \label{eq:patch_embedding}
\end{equation}
where $W_p \in \mathbb{R}^{d_{\mathrm{model}}\times pn}$ and $b_p \in \mathbb{R}^{d_{\mathrm{model}}}$ are learnable parameters.

This tokenization performs two functions simultaneously. It aggregates short local temporal patterns before global sequence modeling, and it reduces the effective token length from $P$ to $P/p$. As a result, the self-attention stage operates on a shorter sequence, which is useful when the objective is to learn a compact temporal lifting module rather than to maximize raw sequence length.

\subsection{Transformer Encoder and Lifted State Construction}

The patch embeddings are augmented with sinusoidal positional encoding and processed by a Transformer encoder:
\begin{equation}
    \tilde{e}_k^{(i)} = e_k^{(i)} + s^{(i)},
\end{equation}
where $s^{(i)}$ denotes the fixed positional encoding associated with token index $i$. Let
\begin{equation}
    H_k = \mathrm{Transformer}(\tilde{e}_k^{(1)},\ldots,\tilde{e}_k^{(N)})
\end{equation}
be the encoded token sequence. The implementation then applies a pooling-based readout operator,
\begin{equation}
    r_k = \mathrm{Pool}(H_k),
\end{equation}
where the pooling rule can be mean, max, or last-token selection, followed by a linear projection
\begin{equation}
    h_k = W_h r_k + b_h \in \mathbb{R}^{d-n}.
    \label{eq:history_feature}
\end{equation}
The lifted state is constructed as
\begin{equation}
    z_k = [x_k;h_k] \in \mathbb{R}^{d},
    \label{eq:lifted_state_form}
\end{equation}
where $x_k$ is the final state in the history window and $h_k$ is the learned temporal feature. The first $n$ coordinates coincide with the physical state, and the fixed projection $x_k=Cz_k$ with $C=[I_n\;0]$ recovers the physical state exactly.

\subsection{Controlled Linear Latent Dynamics and Exact State Recovery}

Once the lifted state has been computed, the latent dynamics are modeled by the linear controlled system
\begin{equation}
    z_{k+1} = A z_k + B u_k,
    \label{eq:learned_dynamics}
\end{equation}
where $A \in \mathbb{R}^{d\times d}$ and $B \in \mathbb{R}^{d\times m}$ are identified matrices. The fixed projection $x_k=Cz_k$ with $C=[I_n\;0]$ recovers the physical state exactly, so no learned nonlinear decoder is needed.

\subsection{Hybrid Training With Encoder Pretraining and EDMD Regression}

The training pipeline combines neural temporal lifting with closed-form linear regression in the learned latent space.

\textbf{Stage 0: Initial EDMD estimate.} With the encoder at its current initialization, all training windows are mapped to latent states and an initial estimate of $(A,B)$ is obtained by regularized least squares.

\textbf{Stage 1: Pretraining.} Let the shifted next-window be
\begin{equation}
    X_k^{+} = [x_{k-P+2}^\top,\,x_{k-P+3}^\top,\,\ldots,\,x_{k+1}^\top]^\top,
\end{equation}
and define the true next latent state by
\begin{equation}
    z_k^{+} = g_\theta(X_k^{+}).
\end{equation}
The predicted next latent and physical states are
\begin{equation}
    \hat{z}_{k+1} = Az_k + Bu_k,
    \qquad
    \hat{x}_{k+1} = C\hat{z}_{k+1}.
\end{equation}
The two principal pretraining losses are the physical-state prediction loss and the latent prediction loss,
\begin{align}
    \mathcal{L}_{x} &= \|\hat{x}_{k+1}-x_{k+1}\|_2^2,
    \\
    \mathcal{L}_{z} &= \|\hat{z}_{k+1}-z_k^{+}\|_2^2.
\end{align}
In the implementation, training also includes a small auxiliary encoder-side regularization term. Since this term is not central to the modeling contribution and does not change the overall three-stage training structure, we treat it only as a minor training detail. The overall pretraining objective is therefore written as
\begin{equation}
    \mathcal{L}_{\mathrm{pre}} = \lambda_x\mathcal{L}_x + \lambda_z\mathcal{L}_z + \lambda_a\mathcal{L}_{\mathrm{aux}}.
    \label{eq:pretraining_loss}
\end{equation}

\textbf{Stage 2: Final EDMD refit.} After the encoder has been pretrained, the latent trajectories are recomputed and the linear dynamics are refit in closed form. Let
\begin{equation}
    Z=[z_1,\ldots,z_M]^\top,
    \qquad
    Z^{+}=[z_1^{+},\ldots,z_M^{+}]^\top,
    \qquad
    U=[u_1,\ldots,u_M]^\top,
\end{equation}
and define the augmented regressor matrix $\Omega=[Z\;U]$. The ridge-regularized EDMD estimate is
\begin{equation}
\begin{aligned}
    \Theta &= (Z^{+})^\top \Omega
    \left(\Omega^\top\Omega + \lambda I\right)^{-1},
\end{aligned}
    \label{eq:edmd_solution}
\end{equation}
where $\Theta \in \mathbb{R}^{d\times(d+m)}$ is partitioned as
\begin{equation}
    \Theta = [A\;B].
\end{equation}
The codebase further supports optional singular-value clipping of $A$ after regression. Since this step is configuration-dependent, we treat it as an implementation safeguard rather than a universal theoretical guarantee.

\subsection{Autoregressive Multi-Step Prediction}

For multi-step rollout, the model is initialized with a real history window and a future control sequence. At each prediction step, the current window is re-encoded, propagated once by \eqref{eq:learned_dynamics}, projected back to the physical state, and then updated by appending the newly predicted state while discarding the oldest state in the window. Thus, multi-step prediction is not obtained by propagating a single latent state indefinitely; it is an autoregressive process that repeatedly re-evaluates the history-based lifting map.

\section{Control-Oriented Interface}
\label{sec:control}

Although the present paper is centered on model learning rather than completed controller validation, the learned predictor is intended to remain compatible with downstream tracking control. This section therefore summarizes the identified model interface that a later controller would use and clarifies what can, and cannot, be claimed at the current stage.

Assume that a reference state trajectory $x_k^{\mathrm{ref}}$ and a corresponding reference input sequence $u_k^{\mathrm{ref}}$ are available. From the reference history window
\begin{equation}
    X_k^{\mathrm{ref}} = [(x_{k-P+1}^{\mathrm{ref}})^\top,\ldots,(x_k^{\mathrm{ref}})^\top]^\top,
\end{equation}
a lifted reference can be formed as
\begin{equation}
    z_k^{\mathrm{ref}} = g_\theta(X_k^{\mathrm{ref}}).
    \label{eq:reference_lifting}
\end{equation}
The lifted tracking error is then
\begin{equation}
    e_k = z_k - z_k^{\mathrm{ref}}.
    \label{eq:tracking_error}
\end{equation}
Under the standard assumption that the learned pair $(A,B)$ is used as the prediction model in a local linear tracking design, one may introduce the control deviation $\delta u_k=u_k-u_k^{\mathrm{ref}}$ and consider the quadratic objective
\begin{equation}
    J = \sum_{k=0}^{\infty}
    \left(e_k^\top Q e_k + \delta u_k^\top R\,\delta u_k\right),
    \label{eq:lqr_objective}
\end{equation}
where $Q\succeq 0$ and $R\succ 0$. The corresponding linear feedback law is written as
\begin{equation}
    \delta u_k = -K_c e_k,
    \label{eq:lqr_control}
\end{equation}
where $K_c$ is the feedback gain computed for the pair $(A,B)$. The notation $K_c$ is used deliberately to avoid conflict with the regression matrix in \eqref{eq:edmd_solution}.

Because physical-state recovery is exact through \eqref{eq:decoder_projection}, the state tracking error satisfies
\begin{equation}
    x_k - x_k^{\mathrm{ref}} = C e_k.
\end{equation}
This makes the model convenient for later controller integration: physical-state tracking metrics can be evaluated directly without introducing a separate learned observation map. In the present manuscript, this section is included to document the control-oriented interface of the learned predictor. The intended comparison is to keep the controller structure fixed and vary only the learned Koopman model, so that differences in tracking behavior can be attributed primarily to modeling quality.

\section{Practical Analysis}
\label{sec:analysis}

This section collects several practical observations that help clarify why the proposed structure is useful and what its limitations are.

\paragraph{Finite-history lifting as compact state augmentation.}
Delay coordinates and history stacks have long been used to enrich the representation of nonlinear dynamics. The proposed model adopts the same intuition but uses a more compact parameterization. Instead of taking the full history stack as the lifted state, it learns a temporal summary $h_k$ from the history window and appends this summary to the current physical state. In this sense, the model can be viewed as a compact learned state augmentation. The current state remains explicit, while recent history is compressed into a smaller auxiliary coordinate block.

\paragraph{Temporal patching as structured compression.}
The use of temporal patches is not only an architectural convenience; it also changes the computational and statistical role of the encoder. A window of length $P$ is reduced to $N=P/p$ tokens, so the sequence length presented to the self-attention block is shorter than the raw time-step sequence. This yields a lower attention cost of order $\mathcal{O}(N^2)$ instead of $\mathcal{O}(P^2)$ with respect to sequence length, while each token already summarizes a short local trajectory segment. For a small-data controlled-system setting, this structured compression is useful because it limits model size without removing temporal order.

\paragraph{State-embedded structure.}
The lifted vector $z_k=[x_k;h_k]$ includes the physical state in its first $n$ coordinates, so the fixed projection $x_k=Cz_k$ recovers the physical state exactly. This is a straightforward structural choice that keeps the model-controller interface simple and interpretable — it does not by itself improve the linear Koopman fit.

\paragraph{Implications for rollout behavior.}
Open-loop rollout accuracy depends on both the quality of the latent regression and the sensitivity of the encoder under repeated re-encoding. One-step propagation in latent space is linear, but multi-step prediction repeatedly feeds predicted states back into the history window. Consequently, long-horizon error is influenced by the operator matrices $(A,B)$, the encoder map $g_\theta$, and the interaction between them. The ridge-regularized EDMD fit reduces ill-conditioning in the latent regression, and optional singular-value clipping can further moderate operator growth when enabled, but the present work does not claim a formal stability guarantee.

\paragraph{Current limitation of the lifting input.}
The present encoder uses state history but not input history. This matches the training and data-construction procedure used here and is sufficient to test the main modeling hypothesis of the paper, namely that temporally informed lifting improves Koopman model quality. Nevertheless, for systems in which recent input history carries essential information that is not recoverable from state history alone, extending the encoder to joint state-input histories is a natural next step.

\section{Experiments}
\label{sec:experiments}

The experimental study consists of completed trajectory-prediction comparison and ablation analysis, followed by a brief description of the intended trajectory-tracking validation setting. The reported evidence in this paper is therefore limited to model-identification quality under open-loop evaluation.

\subsection{Experimental Setup}

\subsubsection{Platform-Specific Datasets}

The present study uses two platform-specific controlled-system datasets with complementary roles.

\textbf{Platform A: manipulator dataset.} This dataset corresponds to a 6-state, 3-input controlled mechanical system. The available split contains 15 training trajectories, 4 validation trajectories, and 1 test trajectory, with 74,655, 19,909, and 4,976 recorded time steps, respectively.

\textbf{Platform B: two-state dataset.} This dataset corresponds to a second controlled setting with state dimension $n=2$ and control dimension $m=2$. The available split contains 49 training trajectories, 21 validation trajectories, and 1 test trajectory, with 24,150, 10,351, and 500 recorded time steps, respectively.

In both datasets, valid samples are constructed from history windows that do not cross trajectory boundaries. State and control normalization statistics are computed on the training split and then reused for validation and test data. The two datasets do not yet form a perfectly uniform benchmark; instead, they play complementary roles in the present paper. Platform A provides the most complete cross-model rollout comparison, whereas Platform B provides the broader hyperparameter sweep used for ablation.

\subsubsection{Compared Methods}

The comparative study is restricted to Koopman model families supported by the present study and evaluated under a common testing pipeline:
\begin{itemize}
    \item \textbf{Traditional EDMD}, which uses hand-crafted lifted observables.
    \item \textbf{MLP-Koopman}, which learns an instantaneous-state lifting map with an MLP encoder.
    \item \textbf{Transformer-Koopman}, the proposed temporal lifting model based on history windows and patch-based Transformer encoding.
\end{itemize}

\subsubsection{Evaluation Protocol}

For completed prediction experiments, all methods are evaluated under the same open-loop autoregressive rollout protocol. Each model is initialized with a real history window from the test trajectory and recursively feeds back its own predicted state. The main metrics are rollout RMSE and MAE over the predicted trajectory.

For the ablation study, all reported configurations on Platform B were trained and tested under the same pipeline. The resulting normalized one-step test RMSE values therefore provide a controlled comparison of the effect of architectural choices within that dataset.

\subsection{Trajectory Prediction Accuracy Comparison}

The most complete cross-model rollout comparison is currently available on Platform A. Table~\ref{tab:prediction_platform_a} reports the trajectory-level results.

\begin{table}[!t]
\renewcommand{\arraystretch}{1.2}
\caption{Open-loop trajectory prediction results on Platform A}
\label{tab:prediction_platform_a}
\centering
\begin{tabular}{lcc}
\hline\hline
Method & RMSE & MAE \\
\hline
Traditional EDMD & 2.6206 & 1.9544 \\
MLP-Koopman & 0.5074 & 0.2511 \\
Transformer-Koopman & \textbf{0.1768} & \textbf{0.0917} \\
\hline\hline
\end{tabular}
\end{table}

The results in Table~\ref{tab:prediction_platform_a} show a clear pattern. Moving from hand-crafted observables to learned liftings already reduces rollout error substantially, as seen in the gap between traditional EDMD and MLP-Koopman. Introducing finite-history temporal lifting yields a further reduction, indicating that the improvement is not only due to using a neural encoder, but also to how the lifting map is constructed. In the present setting, explicitly encoding short trajectory context produces a materially more accurate Koopman model than instantaneous-state lifting alone.

These results should still be interpreted with the correct scope. Platform A is used as the main cross-model case because it is the dataset for which the available evidence provides the most complete like-for-like rollout comparison. Even so, the result is directly relevant to the paper's modeling objective, since recursive rollout accuracy is the most practical test of whether the learned lifted dynamics are useful.

\subsection{Ablation Study}

The ablation study is conducted on Platform B using the available configuration sweep of the proposed model. The best-performing configuration uses history length $P=6$, patch length $p=2$, latent dimension $d=8$, Transformer dimension $d_{\mathrm{model}}=16$, three Transformer layers, two attention heads, and EDMD regularization $\lambda=10^{-4}$.

Table~\ref{tab:ablation_platform_b} reports representative ablation results.

\begin{table}[!t]
\renewcommand{\arraystretch}{1.2}
\caption{Representative ablation results on Platform B}
\label{tab:ablation_platform_b}
\centering
\begin{tabular}{>{\raggedright\arraybackslash}p{0.62\columnwidth}c}
\hline\hline
Variant & Test RMSE \\
\hline
Best setting ($P=6$, $p=2$, 3 layers, 2 heads) & \textbf{0.01408} \\
Reduce history to $P=4$ & 0.01525 \\
Reduce history to $P=2$ & 0.01710 \\
Degenerate patching ($P=2$, $p=1$) & 0.03362 \\
Use 1 layer and 1 head & 0.01513 \\
Use larger latent/model size ($d=16$, $d_{\mathrm{model}}=32$) & 0.01481 \\
Increase regularization to $\lambda=1$ & 0.01591 \\
Increase regularization to $\lambda=10$ & 0.02153 \\
\hline\hline
\end{tabular}
\\[2pt]
{\footnotesize Normalized one-step test RMSE; lower is better.}
\end{table}

Several trends are consistent across these ablations. First, reducing the history length degrades performance, which supports the premise that short-term temporal information is useful for constructing the lifting map. Second, the large error increase in the degenerate short-history setting shows that the gain is not explained by parameter count alone; the patch structure itself matters. Third, increasing latent and embedding size does not automatically improve performance, suggesting that compact temporal liftings can already be sufficient for this dataset. Finally, overly strong EDMD regularization is detrimental, whereas moderate regularization provides a better trade-off.

From a modeling perspective, these results are important because they connect the architecture to the observed gains. The method benefits from three ingredients working together: enough history to capture local dynamics, patch-based compression before attention, and a regularized latent linear fit that avoids over-conditioning the operator regression.

\subsection{Planned Trajectory-Tracking Validation}

The final component of the intended validation pipeline is controller-based trajectory tracking. Since the present study is centered on offline model learning and open-loop evaluation, the manuscript reports the validation protocol rather than completed closed-loop numbers.

The planned comparison is conceptually simple. The controller structure will be kept fixed, and only the learned Koopman model will be changed. In this way, any change in tracking performance can be attributed primarily to model quality rather than to controller redesign. The candidate models for this comparison are Traditional EDMD, MLP-Koopman, and Transformer-Koopman.

The intended evaluation metrics include tracking RMSE, peak tracking error, and control effort or input smoothness. This protocol is consistent with the control-oriented formulation in Section~\ref{sec:control} and provides a direct path toward a later closed-loop study once the controller pipeline is finalized.

\section{Conclusion}
\label{sec:conclusion}

This paper presented Transformer-Koopman, a temporal lifting framework for learning Koopman models of controlled nonlinear systems from finite history windows. The method uses non-overlapping temporal patches and a Transformer encoder to build a compact history-dependent lifting map, propagates the lifted state with linear controlled dynamics, and preserves exact physical-state recovery through a fixed projection. In this way, the resulting model is not only predictive but also structurally suitable for control-oriented evaluation and later controller design.

The completed experiments support the main identification claim of the paper. On the manipulator dataset, the proposed method achieves substantially lower rollout error than the implemented traditional EDMD and MLP-Koopman baselines. On the second dataset, the ablation study shows that temporal context, patch design, and moderate regularization all contribute meaningfully to model quality. Together, these results indicate that a more accurate Koopman predictor can be obtained when the lifting map is learned from temporally structured inputs rather than from the current state alone.

Closed-loop tracking validation is not yet included in the present manuscript. The most direct next step is therefore to compare the learned predictors under a common controller structure and assess whether improved predictive accuracy also translates into improved control performance. A second natural extension is to study richer lifting inputs, including joint state-input histories, within the same control-oriented Koopman framework.

\begin{thebibliography}{99}

\bibitem{koopman1931hamiltonian}
B. O. Koopman, ``Hamiltonian systems and transformation in Hilbert space,'' \emph{Proc. Natl. Acad. Sci. USA}, vol. 17, no. 5, pp. 315--318, 1931.

\bibitem{koopman1932spectra}
B. O. Koopman and J. von Neumann, ``Dynamical systems of continuous spectra,'' \emph{Proc. Natl. Acad. Sci. USA}, vol. 18, no. 3, pp. 255--263, 1932.

\bibitem{mezic2005spectral}
I. Mezi\'c, ``Spectral properties of dynamical systems, model reduction and decompositions,'' \emph{Nonlinear Dyn.}, vol. 41, pp. 309--325, 2005.

\bibitem{ljung1999system}
L. Ljung, \emph{System Identification: Theory for the User}, 2nd ed. Upper Saddle River, NJ, USA: Prentice Hall, 1999.

\bibitem{vanoverschee1994n4sid}
P. Van Overschee and B. De Moor, ``N4SID: Subspace algorithms for the identification of combined deterministic-stochastic systems,'' \emph{Automatica}, vol. 30, no. 1, pp. 75--93, 1994.

\bibitem{nelles2020nonlinear}
O. Nelles, \emph{Nonlinear System Identification: From Classical Approaches to Neural Networks and Fuzzy Models}, 2nd ed. Cham, Switzerland: Springer, 2020.

\bibitem{rowley2009spectral}
C. W. Rowley, I. Mezi\'c, S. Bagheri, P. Schlatter, and D. S. Henningson, ``Spectral analysis of nonlinear flows,'' \emph{J. Fluid Mech.}, vol. 641, pp. 115--127, 2009.

\bibitem{williams2015data}
M. O. Williams, I. G. Kevrekidis, and C. W. Rowley, ``A data--driven approximation of the Koopman operator: Extending dynamic mode decomposition,'' \emph{J. Nonlinear Sci.}, vol. 25, no. 6, pp. 1307--1346, 2015.

\bibitem{williams2015kernel}
M. O. Williams, C. W. Rowley, and I. G. Kevrekidis, ``A kernel-based method for data-driven Koopman spectral analysis,'' \emph{J. Comput. Dyn.}, vol. 2, no. 2, pp. 247--265, 2015.

\bibitem{arbabi2017ergodic}
H. Arbabi and I. Mezi\'c, ``Ergodic theory, dynamic mode decomposition, and computation of spectral properties of the Koopman operator,'' \emph{SIAM J. Appl. Dyn. Syst.}, vol. 16, no. 4, pp. 2096--2126, 2017.

\bibitem{brunton2022modern}
S. L. Brunton, M. Budi\v{s}i\'c, E. Kaiser, and J. N. Kutz, ``Modern Koopman theory for dynamical systems,'' \emph{SIAM Rev.}, vol. 64, no. 2, pp. 229--340, 2022.

\bibitem{otto2021koopman}
S. E. Otto and C. W. Rowley, ``Koopman operators for estimation and control of dynamical systems,'' \emph{Annu. Rev. Control Robot. Auton. Syst.}, vol. 4, pp. 59--87, 2021.

\bibitem{takeishi2017koopman}
N. Takeishi, Y. Kawahara, and T. Yairi, ``Learning Koopman invariant subspaces for dynamic mode decomposition,'' in \emph{Proc. NeurIPS}, 2017, pp. 1130--1140.

\bibitem{yeung2019learning}
E. Yeung, S. Kundu, and N. Hodas, ``Learning deep neural network representations for Koopman operators of nonlinear dynamical systems,'' in \emph{Proc. ACC}, 2019, pp. 4832--4839.

\bibitem{lusch2018deep}
B. Lusch, J. N. Kutz, and S. L. Brunton, ``Deep learning for universal linear embeddings of nonlinear dynamics,'' \emph{Nat. Commun.}, vol. 9, no. 1, p. 4950, 2018.

\bibitem{morton2019deep}
J. Morton, F. D. Witherden, and M. J. Kochenderfer, ``Deep variational Koopman models: Inferring Koopman observations for uncertainty-aware dynamics modeling and control,'' in \emph{Proc. IJCAI}, 2019, pp. 5113--5119.

\bibitem{azencot2020forecasting}
O. Azencot, N. B. Erichson, V. Lin, and M. W. Mahoney, ``Forecasting sequential data using consistent Koopman autoencoders,'' in \emph{Proc. ICML}, 2020, pp. 475--485.

\bibitem{yuan2022variational}
J. Yuan, A. A. Heydari, and Z. Zhang, ``Variational Koopman models: Slow coherent structures and uncertainty quantification,'' \emph{IEEE Trans. Autom. Control}, vol. 67, no. 11, pp. 5905--5920, 2022.

\bibitem{otto2019linearly}
S. E. Otto and C. W. Rowley, ``Linearly-recurrent autoencoder networks for learning dynamics,'' \emph{SIAM J. Appl. Dyn. Syst.}, vol. 18, no. 1, pp. 558--593, 2019.

\bibitem{rank2019recurrent}
M. Rank, A. G. Georgescu, and S. Kirchner, ``Recurrent Koopman operator for data-driven modeling of nonlinear dynamics,'' \emph{IEEE Control Syst. Lett.}, vol. 3, no. 4, pp. 912--917, 2019.

\bibitem{forgione2021dynonet}
M. Forgione and D. Varagnolo, ``DynoNet: A neural network architecture for learning dynamical systems,'' \emph{Int. J. Adapt. Control Signal Process.}, vol. 35, no. 4, pp. 612--629, 2021.

\bibitem{brunton2016sindy}
S. L. Brunton, J. L. Proctor, and J. N. Kutz, ``Discovering governing equations from data by sparse identification of nonlinear dynamical systems,'' \emph{Proc. Natl. Acad. Sci. USA}, vol. 113, no. 15, pp. 3932--3937, 2016.

\bibitem{proctor2016dmdc}
J. L. Proctor, S. L. Brunton, and J. N. Kutz, ``Dynamic mode decomposition with control,'' \emph{SIAM J. Appl. Dyn. Syst.}, vol. 15, no. 1, pp. 142--161, 2016.

\bibitem{proctor2018generalizing}
J. L. Proctor, S. L. Brunton, and J. N. Kutz, ``Generalizing Koopman theory to allow for inputs and control,'' \emph{SIAM J. Appl. Dyn. Syst.}, vol. 17, no. 1, pp. 909--930, 2018.

\bibitem{brunton2016invariant}
S. L. Brunton, B. W. Brunton, J. L. Proctor, and J. N. Kutz, ``Koopman invariant subspaces and finite linear representations of nonlinear dynamical systems for control,'' \emph{PLOS ONE}, vol. 11, no. 2, e0150171, 2016.

\bibitem{korda2018linear}
M. Korda and I. Mezi\'c, ``Linear predictors for nonlinear dynamical systems: Koopman operator meets model predictive control,'' \emph{Automatica}, vol. 93, pp. 149--160, 2018.

\bibitem{han2020deep}
Y. Han, W. Hao, and U. Vaidya, ``Deep learning of Koopman representation for control,'' arXiv:2010.07546, 2020.

\bibitem{shi2022deep}
H. Shi and M. Q.-H. Meng, ``Deep Koopman operator with control for nonlinear systems,'' \emph{IEEE Robot. Autom. Lett.}, vol. 7, no. 3, pp. 7700--7707, 2022.

\bibitem{strasser2026overview}
R. Str\"asser, K. Worthmann, I. Mezi\'c, J. Berberich, M. Schaller, and F. Allg\"ower, ``An overview of Koopman-based control: From error bounds to closed-loop guarantees,'' \emph{Annu. Rev. Control}, vol. 61, art. 101035, 2026.

\bibitem{duran2025control}
J. F. Dur\'an-Sig\"uenza, L. I. Minchala, L. E. Garza-Casta\~n\'on, and H. Zhang, ``Control based on the Koopman operator: A comprehensive review,'' \emph{J. Franklin Inst.}, vol. 362, no. 18, art. 108256, 2025.

\bibitem{klus2018data}
S. Klus, P. Koltai, and C. Sch\"utte, ``Data-driven model reduction and transfer operator approximation,'' \emph{J. Nonlinear Sci.}, vol. 28, no. 3, pp. 985--1010, 2018.

\bibitem{vaswani2017attention}
A. Vaswani \emph{et al.}, ``Attention is all you need,'' in \emph{Proc. NeurIPS}, 2017, pp. 5998--6008.

\bibitem{zhou2021informer}
H. Zhou \emph{et al.}, ``Informer: Beyond efficient Transformer for long sequence time-series forecasting,'' in \emph{Proc. AAAI}, 2021, pp. 11106--11115.

\bibitem{wu2021autoformer}
H. Wu, J. Xu, J. Wang, and M. Long, ``Autoformer: Decomposition Transformers with auto-correlation for long-term series forecasting,'' in \emph{Proc. NeurIPS}, 2021, pp. 22419--22430.

\bibitem{zhou2022fedformer}
T. Zhou \emph{et al.}, ``FEDformer: Frequency enhanced decomposed Transformer for long-term series forecasting,'' in \emph{Proc. ICML}, 2022, pp. 27268--27286.

\bibitem{nie2023time}
Y. Nie, N. H. Nguyen, P. Sinthong, and J. Kalagnanam, ``A time series is worth 64 words: Long-term forecasting with Transformers,'' in \emph{Proc. ICLR}, 2023.

\bibitem{zeng2023transformers}
A. Zeng, M. Chen, L. Zhang, and Q. Xu, ``Are Transformers effective for time series forecasting?,'' in \emph{Proc. AAAI}, 2023, pp. 11121--11128.

\bibitem{liu2024itransformer}
Y. Liu \emph{et al.}, ``iTransformer: Inverted Transformers are effective for time series forecasting,'' in \emph{Proc. ICLR}, 2024.

\bibitem{forootani2026compact}
A. Forootani and M. Khosravi, ``Compact transformer variants for synthetic time series forecasting,'' \emph{Neurocomputing}, vol. 677, art. 133140, 2026.

\bibitem{forootani2025synthetic}
A. Forootani and M. Khosravi, ``Synthetic time series forecasting with Transformer architectures: Extensive simulation benchmarks,'' arXiv:2505.20048, 2025.

\bibitem{forootani2026deepkoopformer}
A. Forootani, M. Khosravi, and M. Barati, ``DeepKoopFormer: A Koopman enhanced Transformer based architecture for time series forecasting,'' \emph{npj Artif. Intell.}, vol. 2, art. 35, 2026.

\bibitem{forootani2026spectral}
A. Forootani and R. Iervolino, ``Learnable Koopman-enhanced Transformer-based time series forecasting with spectral control,'' arXiv:2602.02592, 2026.

\bibitem{wang2020linformer}
S. Wang, B. Z. Li, M. Khabsa, H. Fang, and H. Ma, ``Linformer: Self-attention with linear complexity,'' arXiv:2006.04768, 2020.

\bibitem{katharopoulos2020transformers}
A. Katharopoulos, A. Vyas, N. Pappas, and F. Fleuret, ``Transformers are RNNs: Fast autoregressive Transformers with linear attention,'' in \emph{Proc. ICML}, 2020, pp. 5156--5165.

\bibitem{choromanski2021performer}
K. M. Choromanski \emph{et al.}, ``Rethinking attention with Performers,'' in \emph{Proc. ICLR}, 2021.

\end{thebibliography}

\end{document}
